\documentclass{amsart}

\usepackage{fullpage}
\usepackage{setspace}
\onehalfspacing

\newtheorem{claim}{Claim}
\newtheorem{lemma}{Lemma}
\newtheorem{conjecture}{Conjecture}
\newtheorem{theorem}{Theorem}

\newcommand{\N}{\mathbb N}
\newcommand{\Pp}{\mathcal P}

\begin{document}

All logarithms are natural.

Define the von Mangoldt function \(\Lambda\) by
\[
\Lambda(q)=
\begin{cases}
\log p, & q=p^k \text{ for some prime } p \text{ and integer } k\geq 1,\\
0, & \text{otherwise.}
\end{cases}
\]

For \(N,M\in\mathbb Z_{>0}\), define \(v(N,M)\) as follows.  If \(N>1\), \(q\) is a prime-power divisor of \(N\), and
\[
M=\frac{N}{q},
\]
then
\[
v(N,M)=\frac{\Lambda(q)}{N\log^2 N}.
\]
In all other cases, set
\[
v(N,M)=0.
\]

Define \(w(N,M)\) by modifying \(v(N,M)\) only in the following cases.  For every prime \(p\) and every integer \(k\geq 2\), set
\[
w(p^k,1)=0
\]
and
\[
w(p^k,p^{k-1})
=
v(p^k,p^{k-1})+v(p^k,1).
\]
For all other pairs \((N,M)\), set
\[
w(N,M)=v(N,M).
\]

For \(N\geq 1\), define
\[
\operatorname{Out}(N)=\sum_{M\geq 1} w(N,M)
\]
and
\[
\operatorname{In}(N)=\sum_{K\geq 1} w(K,N).
\]

\begin{claim}
For every \(N>1\),
\[
\operatorname{Out}(N)=\frac{1}{N\log N}.
\]
\end{claim}

\begin{proof}
Fix \(N>1\).  The passage from \(v\) to \(w\) only moves mass within a fixed row \(N\).  If \(N\) is not of the form \(p^k\) with \(p\) prime and \(k\geq 2\), then
\[
w(N,M)=v(N,M)
\]
for every \(M\).  If \(N=p^k\) with \(k\geq 2\), then the only changed entries in row \(N\) are \(M=1\) and \(M=p^{k-1}\), and
\[
w(p^k,1)+w(p^k,p^{k-1})
=
0+\bigl(v(p^k,p^{k-1})+v(p^k,1)\bigr).
\]
Thus in every case,
\[
\sum_{M\geq 1}w(N,M)=\sum_{M\geq 1}v(N,M).
\]

Now compute the row sum of \(v\).  By definition,
\[
\sum_{M\geq 1}v(N,M)
=
\sum_{\substack{q\mid N\\ q\text{ prime power}}}
\frac{\Lambda(q)}{N\log^2 N}.
\]
Since \(\Lambda(q)=0\) unless \(q\) is a prime power,
\[
\sum_{M\geq 1}v(N,M)
=
\frac{1}{N\log^2 N}\sum_{q\mid N}\Lambda(q).
\]
If
\[
N=\prod_{i=1}^r p_i^{a_i},
\]
then the prime-power divisors contributing to \(\sum_{q\mid N}\Lambda(q)\) are
\[
p_i,p_i^2,\dots,p_i^{a_i},
\]
and each contributes \(\log p_i\).  Therefore
\[
\sum_{q\mid N}\Lambda(q)
=
\sum_{i=1}^r\sum_{j=1}^{a_i}\log p_i
=
\sum_{i=1}^r a_i\log p_i
=
\log N.
\]
Hence
\[
\operatorname{Out}(N)
=
\sum_{M\geq 1}w(N,M)
=
\sum_{M\geq 1}v(N,M)
=
\frac{\log N}{N\log^2 N}
=
\frac{1}{N\log N}.
\]
\end{proof}

\begin{lemma}\label{lem:analytic-estimate}
For \(s>1\), define
\[
A(s)=\sum_{q\geq 2}\frac{\Lambda(q)}{q^s}.
\]
Then, for every prime \(p\),
\[
A(s)+\frac{\log p}{p^s}
\leq
\frac{1}{s-1}.
\]
\end{lemma}

\begin{proof}
For \(s>1\),
\[
A(s)=-\frac{\zeta'(s)}{\zeta(s)},
\]
where
\[
\zeta(s)=\sum_{n=1}^{\infty}\frac{1}{n^s}.
\]
Use the standard estimates
\[
-\zeta'(s)
=
\sum_{n=2}^{\infty}\frac{\log n}{n^s}
\leq
\int_1^\infty \frac{\log x}{x^s}\,dx
=
\frac{1}{(s-1)^2}
\]
and
\[
\zeta(s)\geq \frac{1}{s-1}+\frac12.
\]
The second estimate follows from the convexity of \(x^{-s}\): for each integer \(n\geq 1\),
\[
\int_n^{n+1} x^{-s}\,dx
\leq
\frac{n^{-s}+(n+1)^{-s}}{2},
\]
and summing over \(n\geq 1\) gives
\[
\frac{1}{s-1}
=
\int_1^\infty x^{-s}\,dx
\leq
\zeta(s)-\frac12.
\]

Therefore
\[
A(s)
=
-\frac{\zeta'(s)}{\zeta(s)}
\leq
\frac{\frac{1}{(s-1)^2}}{\frac{1}{s-1}+\frac12}
=
\frac{2}{s^2-1}.
\]
Also, for every prime \(p\),
\[
\frac{\log p}{p^s}
\leq
\max_{x>1}\frac{\log x}{x^s}
=
\frac{1}{es}
\leq
\frac{1}{s+1}.
\]
Hence
\[
A(s)+\frac{\log p}{p^s}
\leq
\frac{2}{s^2-1}+\frac{1}{s+1}
=
\frac{1}{s-1}.
\]
\end{proof}

\begin{claim}
For every \(N>1\),
\[
\operatorname{Out}(N)\geq \operatorname{In}(N).
\]
\end{claim}

\begin{proof}
Fix \(N>1\), and write
\[
L=\log N.
\]
Then \(L>0\).

A nonzero \(v\)-flow into \(N\) comes from a source \(Nq\), where \(q\geq 2\).  For every \(q\geq 2\),
\[
v(Nq,N)=\frac{\Lambda(q)}{Nq\log^2(Nq)}.
\]
Thus the \(v\)-contribution to the inflow into \(N\) is
\[
\sum_{q\geq 2}\frac{\Lambda(q)}{Nq\log^2(Nq)}.
\]

The modification from \(v\) to \(w\) can add extra inflow into \(N>1\) only when \(N\) is a prime power.  If \(N\) is not a prime power, then
\[
N\operatorname{In}(N)
=
\sum_{q\geq 2}
\frac{\Lambda(q)}{q(L+\log q)^2}.
\]
If
\[
N=p^a
\]
with \(p\) prime and \(a\geq 1\), then the extra inflow into \(N\) comes from the modified flow
\[
p^{a+1}\longrightarrow p^a.
\]
The extra amount is
\[
v(p^{a+1},1)
=
\frac{\log p}{p^{a+1}\log^2(p^{a+1})}.
\]
Multiplying by \(N=p^a\), this extra contribution becomes
\[
N v(p^{a+1},1)
=
\frac{\log p}{p(L+\log p)^2}.
\]
Therefore, when \(N=p^a\),
\[
N\operatorname{In}(N)
=
\sum_{q\geq 2}
\frac{\Lambda(q)}{q(L+\log q)^2}
+
\frac{\log p}{p(L+\log p)^2}.
\]

Use the identity
\[
\frac{1}{Y^2}
=
\int_0^\infty t e^{-tY}\,dt
\qquad (Y>0).
\]
If \(N=p^a\), then
\[
\begin{aligned}
N\operatorname{In}(N)
&=
\int_0^\infty
t e^{-Lt}
\left(
\sum_{q\geq 2}\frac{\Lambda(q)}{q^{1+t}}
+
\frac{\log p}{p^{1+t}}
\right)
\,dt.
\end{aligned}
\]
By Lemma \ref{lem:analytic-estimate}, with \(s=1+t\),
\[
\sum_{q\geq 2}\frac{\Lambda(q)}{q^{1+t}}
+
\frac{\log p}{p^{1+t}}
\leq
\frac{1}{t}.
\]
Thus
\[
N\operatorname{In}(N)
\leq
\int_0^\infty t e^{-Lt}\cdot \frac{1}{t}\,dt
=
\int_0^\infty e^{-Lt}\,dt
=
\frac{1}{L}.
\]
Hence
\[
\operatorname{In}(N)\leq \frac{1}{N\log N}.
\]

If \(N\) is not a prime power, the same argument applies without the extra term, using
\[
\sum_{q\geq 2}\frac{\Lambda(q)}{q^{1+t}}
\leq
\frac{1}{t}.
\]
Therefore again
\[
\operatorname{In}(N)\leq \frac{1}{N\log N}.
\]

By Claim \(1\),
\[
\operatorname{Out}(N)=\frac{1}{N\log N}.
\]
Therefore
\[
\operatorname{Out}(N)\geq \operatorname{In}(N).
\]
\end{proof}

\begin{claim}
For every \(N>1\),
\[
w(N,1)>0
\quad\Longleftrightarrow\quad
N \text{ is prime}.
\]
\end{claim}

\begin{proof}
A nonzero \(v\)-flow from \(N\) to \(1\) can occur only if
\[
1=\frac{N}{q},
\]
so
\[
q=N.
\]
Thus, for \(N>1\),
\[
v(N,1)=\frac{\Lambda(N)}{N\log^2 N},
\]
and this is nonzero exactly when \(N\) is a prime power.

If \(N\) is not a prime power, then
\[
\Lambda(N)=0,
\]
so
\[
v(N,1)=0.
\]
No exceptional modification applies, and therefore
\[
w(N,1)=0.
\]

If \(N=p^k\) with \(p\) prime and \(k\geq 2\), then \(v(N,1)>0\), but this is exactly one of the exceptional cases in the definition of \(w\).  Hence
\[
w(p^k,1)=0.
\]

Finally, if \(N=p\) is prime, then no exceptional modification applies, so
\[
w(p,1)=v(p,1)
=
\frac{\Lambda(p)}{p\log^2 p}
=
\frac{\log p}{p\log^2 p}
=
\frac{1}{p\log p}
>
0.
\]
Therefore
\[
w(N,1)>0
\quad\Longleftrightarrow\quad
N \text{ is prime}.
\]
\end{proof}

Let
\[
\N_{\geq 2}:=\{2,3,4,\dots\},
\]
and let \(\Pp\) denote the set of primes.  A set
\[
A\subseteq \N_{\geq 2}
\]
is called \emph{primitive} if no two distinct elements of \(A\) divide one another.

A \emph{divisibility flow} is a nonnegative function \(W(m,n)\) on pairs of positive integers such that
\[
W(m,n)=0
\]
unless
\[
n\mid m
\qquad\text{and}\qquad
n<m.
\]

\begin{conjecture}\label{conj:flow}
There exists a divisibility flow \(w\) such that, for every \(r\geq 2\),
\[
\operatorname{Out}(r)\geq \frac{1}{r\log r},
\]
with equality whenever \(r\) is prime;
\[
\operatorname{Out}(r)\geq \operatorname{In}(r);
\]
and the only nonzero flows into \(1\) are those coming from primes:
\[
w(m,1)=0
\qquad
\text{whenever } m\geq 2 \text{ is composite}.
\]
\end{conjecture}

\begin{proof}
The function \(w\) defined above is nonnegative and is zero unless \(M\mid N\) and \(M<N\).  Thus it is a divisibility flow.

Claim \(1\) gives
\[
\operatorname{Out}(r)=\frac{1}{r\log r}
\]
for every \(r\geq 2\).  In particular,
\[
\operatorname{Out}(r)\geq \frac{1}{r\log r},
\]
with equality whenever \(r\) is prime.

Claim \(2\) gives
\[
\operatorname{Out}(r)\geq \operatorname{In}(r)
\]
for every \(r\geq 2\).

Claim \(3\) gives
\[
w(m,1)=0
\]
for every composite \(m\geq 2\).  Hence all three asserted properties hold.
\end{proof}

\begin{theorem}[Main result]
For every primitive set \(A\subseteq \N_{\geq 2}\),
\[
\sum_{a\in A}\frac{1}{a\log a}
\leq
\sum_{p\in\Pp}\frac{1}{p\log p}.
\]
Consequently,
\[
\sum_{n\in A}\frac{1}{n\log n}
\]
is maximized over primitive sets \(A\subseteq \N_{\geq 2}\) by taking \(A=\Pp\).
\end{theorem}

\begin{proof}
First suppose that \(A\) is finite.  Define
\[
\Omega
:=
\{d\geq 2:\ d\mid a \text{ for some } a\in A\}.
\]
Then \(\Omega\) is finite, \(A\subseteq \Omega\), and \(\Omega\) is downward closed above \(1\): if \(d\in\Omega\), \(e\mid d\), and \(e\geq 2\), then \(e\in\Omega\).

For this finite set \(\Omega\), define
\[
\operatorname{Out}(\Omega)
:=
\sum_{\substack{m\in\Omega,\ n\notin\Omega\\ n\mid m,\ n<m}}
w(m,n)
\]
and
\[
\operatorname{In}(\Omega)
:=
\sum_{\substack{m\notin\Omega,\ n\in\Omega\\ n\mid m,\ n<m}}
w(m,n).
\]
The finite divergence identity is
\[
\sum_{r\in\Omega}
\bigl(\operatorname{Out}(r)-\operatorname{In}(r)\bigr)
=
\operatorname{Out}(\Omega)-\operatorname{In}(\Omega).
\]
Indeed, after expanding the left-hand side, every term \(w(m,n)\) with both \(m,n\in\Omega\) occurs once with a positive sign and once with a negative sign.  Terms leaving \(\Omega\) occur once with a positive sign, and terms entering \(\Omega\) occur once with a negative sign.

We first compute the possible nonzero terms in \(\operatorname{Out}(\Omega)\).  Let \(m\in\Omega\), and suppose \(n\notin\Omega\), \(n\mid m\), and \(n<m\).  If \(n\geq 2\), then the downward closure of \(\Omega\) gives \(n\in\Omega\), a contradiction.  Hence \(n=1\).  By Conjecture \ref{conj:flow}, the only nonzero flows into \(1\) come from primes.  Therefore
\[
\operatorname{Out}(\Omega)
=
\sum_{\substack{p\in\Omega\\ p\text{ prime}}}
w(p,1).
\]
If \(p\) is prime, then its only proper divisor is \(1\), so
\[
w(p,1)=\operatorname{Out}(p).
\]
By Conjecture \ref{conj:flow},
\[
\operatorname{Out}(p)=\frac{1}{p\log p}.
\]
Thus
\[
\operatorname{Out}(\Omega)
=
\sum_{\substack{p\in\Omega\\ p\text{ prime}}}
\frac{1}{p\log p}
\leq
\sum_{p\in\Pp}\frac{1}{p\log p}.
\]

Next, the finite divergence identity gives
\[
\operatorname{Out}(\Omega)
=
\operatorname{In}(\Omega)
+
\sum_{r\in\Omega}
\bigl(\operatorname{Out}(r)-\operatorname{In}(r)\bigr).
\]
By Conjecture \ref{conj:flow},
\[
\operatorname{Out}(r)-\operatorname{In}(r)\geq 0
\]
for every \(r\geq 2\).  Since \(A\subseteq\Omega\),
\[
\operatorname{Out}(\Omega)
\geq
\operatorname{In}(\Omega)
+
\sum_{a\in A}
\bigl(\operatorname{Out}(a)-\operatorname{In}(a)\bigr).
\]

We claim that every nonzero flow entering a point \(a\in A\) enters \(\Omega\) from outside \(\Omega\).  Such a flow has the form
\[
m\longrightarrow a
\]
with
\[
a\mid m
\qquad\text{and}\qquad
m>a.
\]
If \(m\in\Omega\), then \(m\mid b\) for some \(b\in A\).  Hence
\[
a\mid m\mid b.
\]
Since \(A\) is primitive and \(a,b\in A\), this forces \(a=b\).  Then \(m\mid a\), contradicting \(m>a\).  Therefore \(m\notin\Omega\).

It follows that the full inflow into every \(a\in A\) is counted in \(\operatorname{In}(\Omega)\).  Hence
\[
\operatorname{In}(\Omega)
\geq
\sum_{a\in A}\operatorname{In}(a).
\]
Combining this with the preceding inequality gives
\[
\operatorname{Out}(\Omega)
\geq
\sum_{a\in A}\operatorname{In}(a)
+
\sum_{a\in A}
\bigl(\operatorname{Out}(a)-\operatorname{In}(a)\bigr)
=
\sum_{a\in A}\operatorname{Out}(a).
\]
By Conjecture \ref{conj:flow},
\[
\operatorname{Out}(a)\geq \frac{1}{a\log a}
\]
for every \(a\in A\).  Therefore
\[
\sum_{a\in A}\frac{1}{a\log a}
\leq
\sum_{a\in A}\operatorname{Out}(a)
\leq
\operatorname{Out}(\Omega)
\leq
\sum_{p\in\Pp}\frac{1}{p\log p}.
\]
This proves the desired inequality when \(A\) is finite.

Now let \(A\subseteq\N_{\geq 2}\) be an arbitrary primitive set.  For every finite subset \(A_0\subseteq A\), the finite case gives
\[
\sum_{a\in A_0}\frac{1}{a\log a}
\leq
\sum_{p\in\Pp}\frac{1}{p\log p}.
\]
Taking the supremum over all finite subsets \(A_0\subseteq A\) gives
\[
\sum_{a\in A}\frac{1}{a\log a}
\leq
\sum_{p\in\Pp}\frac{1}{p\log p}.
\]

The set of primes \(\Pp\) is primitive, and for \(A=\Pp\) the left-hand side is equal to the right-hand side.  Hence the maximum over primitive sets is attained by the set of primes.
\end{proof}

\end{document}
